% Akronimy: nazwa symboliczna (dla LaTeXa), nazwa skrócona, nazwa pełna
\newacronym{AMP}{AMP}{Asymmetric Multiprocessing}
\newacronym{SoC}{SoC}{System on a~Chip}
\newacronym{IPC}{IPC}{Inter-Processor Communication}
\newacronym{FPGA}{FPGA}{Field-Programmable Gate Array}
\newacronym{HPS}{HPS}{Hard Processor System}
\newacronym{DTB}{DTB}{Device Tree Blob}
\newacronym{DTS}{DTS}{Device Tree Source}
\newacronym{RT}{RT}{Real-Time}
\newacronym{VPN}{VPN}{Virtual Private Network}
\newacronym{UART}{UART}{Universal Asynchronous Receiver-Transmitter}
\newacronym{SSH}{SSH}{Secure Shell}
\newacronym{SMC}{SMC}{Secure Monitor Call}
\newacronym{MAC}{MAC}{Media Access Control}
\newacronym{IP}{IP}{Internet Protocol}
\newacronym{RAM}{RAM}{Random Access Memory}
\newacronym{CLI}{CLI}{Command Line Interface}
\newacronym{IPI}{IPI}{Inter-Processor Interrupt}

% Wpisy do słownika pojęć (Glossary)
\newglossaryentry{baremetal}{
    name={Bare-metal}, 
    description={Oprogramowanie wykonywane bezpośrednio na sprzęcie, bez pośrednictwa systemu operacyjnego oraz bez mechanizmów zarządzania pamięcią wirtualną.}
}

\newglossaryentry{ringbuffer}{
    name={Bufor kołowy}, 
    description={Struktura danych o~stałym rozmiarze (ang. \textit{ring buffer}), w~której koniec bufora łączy się z~jego początkiem. W~projekcie wykorzystywana do bezkolizyjnej komunikacji IPC.}
}

\newglossaryentry{lwh2f}{
    name={LWH2F}, 
    description={Lightweight HPS-to-FPGA -- lekki mostek komunikacyjny pomiędzy systemem procesora (HPS) a~matrycą programowalną (FPGA).}
}

\newglossaryentry{yocto}{
    name={Yocto Project}, 
    description={Środowisko (tzw. \textit{build system}) służące do tworzenia zoptymalizowanych, dedykowanych dystrybucji systemu Linux dla systemów wbudowanych.}
}

\newglossaryentry{sysfs}{
    name={sysfs}, 
    description={Wirtualny system plików w~systemie Linux, który eksponuje informacje o~urządzeniach i~sterownikach do przestrzeni użytkownika.}
}

\newglossaryentry{tailscale}{
    name={Tailscale}, 
    description={Oprogramowanie tworzące wirtualną sieć prywatną w~topologii kraty (ang. \textit{mesh VPN}) opartą na protokole WireGuard.}
}